% !PRINTED 498
\weberpage{498}
%<<CONT
woraus folgt, daß für keinen anderen Wert von $u^{24}$ als 64 zwei
der Werte $v_\nu$ einander gleich werden. Es tritt diese Gleichheit
also nur für $\omega = i$ und gewisse mit $i$ äquivalente Werte von $\omega$
ein, und daher nicht für einen anderen rein imaginären Wert
von $\omega$.

Der Anfang der Entwickelung nach steigenden Potenzen von
$q$ ergibt nun:
\weq{26}{\begin{aligned}
v_2 - v_5 &= f\left(\frac{\omega - 2\,.\,48}{7}\right) - f\left(\frac{\omega + 2\,.\,48}{7}\right)\\
&= 2\,i\,q^{-\frac{1}{7\,.\,24}} \sin\frac{4\pi}{7} + \cdots,
\end{aligned}}
so daß für einen hinlänglich großen reellen Wert von $-i\omega$ die
linke Seite von (26) positiv imaginär ist. Wenn nun $-i\omega$ auf
reellem Wege von $\infty$ bis 1 geht, so geht, wie wir oben gesehen
haben, $v_2 - v_5$ nicht durch Null. Es muß also auch für $\omega = i$
die Differenz $v_2 - v_5$ positiv imaginär sein, wie in (24) angenommen
 ist.

Wenn man also die Werte (23), (24) in (14) einsetzt, so
ergibt sich für $\omega = i$:
\[
w_0 = 0, \quad w_5 = 0
\]
\weq{27}{\begin{aligned}
w_1 = w_2 = w_3 = w_6 &= \frac{7 + i\sqrt{7}}{2},\\
w_4 &= -\frac{i\sqrt{7}\,(1+i\sqrt{7})^2}{4}.
\end{aligned}}
Danach kann für den Fall $n = 7$ die Resolvente vollständig
gebildet werden. Sie lautet:
\weq{28}{\begin{aligned}
&w^2\left(w - \frac{7+i\sqrt7}{2}\right)^4\left(w + \frac{i\sqrt7\,(1+i\sqrt7)^2}{4}\right)\\
&= \left(u^{12} - \frac{64}{u^{12}}\right)^2,
\end{aligned}}
und nimmt eine einfachere Gestalt an, wenn man
\[
w + \frac{i\sqrt7\,(1+i\sqrt7)^2}{4} = z^2
\]
setzt:
\weq{29}{z\left(z^2 - \frac{i\sqrt7\,(1+i\sqrt7)^2}{4}\right)\left(z^2+i\sqrt7\right)^2 = u^{12} - \frac{64}{u^{12}}.}
